
%(BEGIN_QUESTION)
% Copyright 2006, Tony R. Kuphaldt, released under the Creative Commons Attribution License (v 1.0)
% This means you may do almost anything with this work of mine, so long as you give me proper credit

En person flyter i et svømmebasseng på en oppblåsbar plastpute og holder en kopp drikkevann. Hvis personen ved et uhell heller drikkevannet sitt i bassenget, vil vannstanden i bassenget øke, synke eller forbli den samme? Forklar hvorfor.

\underbar{file i00800}
%(END_QUESTION)





%(BEGIN_ANSWER)

Vannstanden i bassenget forblir den samme, fordi vannet som fortrenges av den oppblåsbare puten er en vannmengde med samme vekt som puten, personen, koppen og vannet i koppen til sammen. Når vannet i koppen helles ut i bassenget, blir lasten på puten mindre, slik at den fortrenger mindre vann, men denne reduserte fortrengningen er {\it nøyaktig lik} vannmengden som ble helt i bassenget. Derfor blir det ingen endring i vannstanden.

%(END_ANSWER)





%(BEGIN_NOTES)


%INDEX% Physics, static fluids: buoyancy

%(END_NOTES)
